\documentclass[10pt,a4paper]{amsart}
\usepackage[margin=19mm]{geometry}
\usepackage{amsmath,amssymb,mathtools}
\usepackage[scr=rsfs,cal=euler]{mathalfa}
\usepackage{xfrac}
\usepackage[unicode,hidelinks,bookmarks=false]{hyperref}
\newcommand{\ie}{i.e.\ }
\newcommand{\cf}{cf.\ }
\newcommand{\resp}{respectively\ }
\newcommand{\Subsection}{\subsection}
\providecommand{\nobreakdash}{\nobreak\hbox{-}}
\setlength{\emergencystretch}{2em}
\multlinegap0pt
\usepackage{mathtools}
\usepackage{amssymb}
\usepackage{xcolor}
\newtheorem{lemma}{Lemma}
\newtheorem{proposition}{Proposition}
\def \N {\mathbb{N}}
\def \Q {\mathcal Q}
\def \R {\mathbb{R}}
\renewcommand{\ge}{\geqslant}
\renewcommand{\geq}{\geqslant}
\renewcommand{\le}{\leqslant}
\renewcommand{\leq}{\leqslant}
\title{Some nonlinear problems\\ for the superposition of fractional operators\\ with Neumann boundary conditions}
\author{Serena Dipierro, Edoardo Proietti Lippi, Caterina Sportelli, Enrico Valdinoci}
\date{Source: arXiv:2404.11091v1 (CC BY 4.0)}
\begin{document}
\maketitle

\noindent\emph{Source and licence.} Some nonlinear problems\\ for the superposition of fractional operators\\ with Neumann boundary conditions. Version arXiv:2404.11091v1 (CC BY 4.0), distributed by arXiv under the Creative Commons Attribution 4.0 International licence, http://creativecommons.org/licenses/by/4.0/, as recorded in the arXiv licence metadata for this exact version and shown on the abstract page https://arxiv.org/abs/2404.11091v1. Full licence notice: \texttt{licenses/arxiv-2404.11091-CC-BY-4.0.txt}.\par\smallskip

\noindent\textit{Editorial scope.} Counted unit: the article's proposition boundedlinking (source0.tex lines 896--907). The article's complete original proof of it is delivered with it (source0.tex lines 909--966, one \texttt{proof} environment of 58 lines). No mathematics is written, repaired or re-derived: the statement and the proof are the article's own lines, in the article's own notation, and no retained line is substituted or re-flowed.  Retained uncounted as the source-local context the proof needs: lines 210--215; lines 217--219; lines 221--223; lines 225--230; lines 232--237; lines 593--604; lines 854--856; lines 877--882; lines 884--889.  Nothing was omitted from any retained span, so the package carries no omission record.  No retained line contains a citation, so the wrapper carries no bibliography and no bibliography entry.  No retained line contains a figure environment, an image-inclusion command or drawing code, so no image is delivered and the package declares no asset dependency.  Editorial formatting only, in addition: an AMS \texttt{amsart} preamble that restates the article's own declarations and macros used by the retained lines, namely the environments \texttt{lemma}, \texttt{proposition} on separate counters, together with the standard packages those lines need.  The retained environments print in this document as Lemma 1, Proposition 1 in order of appearance, whereas the article prints the same items with the numbering of its own full document.  The complete native source of the article as distributed in the arXiv e-print for this exact version is bundled unchanged as \texttt{sources/157788/source0.tex}.
\begin{equation}\label{AR1}
\begin{split}
&\mbox{there exist some constants $a_1,a_2>0$ and $p\in (2, 2^*_{s_\sharp})$ such that} \\
&|f(x,t)|\leq a_1+a_2|t|^{p-1} \quad\mbox{for any } t\in \R \mbox{ and for a.e. } x\in \Omega,\\
\end{split}
\end{equation}

\begin{equation}\label{AR2}
\lim_{t\to 0}\frac{f(x,t)}{t}=0  \quad \mbox{ uniformly for a.e. } x\in \Omega.
\end{equation}

\begin{equation}\label{definizioneF}
F(x,t):=\int_0^t f(x,\tau)\,d\tau,
\end{equation}

\begin{equation}\label{AR3}
\begin{split}
&\mbox{there exist $\vartheta>2$ and $r\geq 0$ such that}\\
&0<\vartheta F(x,t)\leq f(x,t)t  \quad\mbox{for any } t \mbox{ such that } |t|> r \mbox{ and for a.e. } x\in \Omega.
\end{split}
\end{equation}

\begin{equation}\label{AR4}
\begin{split}
&\mbox{there exist some constants~$\widetilde{\vartheta}>2$, $a_3>0$ and a function~$a_4\in L^1(\Omega)$ such that}\\
&F(x,t)\geq a_3|t|^{\widetilde{\vartheta}}-a_4(x).
\end{split}
\end{equation}

\begin{lemma}\label{SV12}
Assume that $f$ satisfies~\eqref{AR1} and~\eqref{AR2}.
Then, for any~$\varepsilon>0$ there exists~$\delta=\delta(\varepsilon)$ such that, for a.e.~$x\in \Omega$ and for any~$t\in \R$,
\begin{eqnarray}&&
|f(x,t)|\leq 2\varepsilon|t|+p\delta(\varepsilon)|t|^{p-1}\nonumber
\\
{\mbox{and }}&&
\label{F1}
|F(x,t)|\leq \varepsilon|t|^2+\delta(\varepsilon)|t|^{p},
\end{eqnarray}
where $F$ is defined as in~\eqref{definizioneF}.
\end{lemma}

\begin{equation}\label{successivi}
\lambda\in [\lambda_i, \lambda_{i+1})\quad\mbox{ for some } i\in\N\setminus\{0\}.
\end{equation}

\begin{equation}\label{funlink}
\begin{split}
\mathcal I(u):&= \frac{\alpha}{2}\int_\Omega |\nabla u|^2 \,dx +\int_{(0, 1)} \frac{c_{N, s}}{4} \iint_{\Q} \frac{|u(x)-u(y)|^2}{|x-y|^{N+2s}} \,dx\,dy\, d\mu(s)\\
&\qquad+\frac12 \int_{\Omega} |u|^2 \, dx -\frac{\lambda}{2} \int_{\Omega} |u|^2 \, dx - \int_\Omega F(x, u) \, dx
\end{split}
\end{equation}

\begin{equation}\label{diff2}
\begin{split}
\langle \mathcal I'(u), v\rangle &=\alpha \int_\Omega \nabla u\cdot\nabla v\, dx +\int_{(0, 1)} \frac{c_{N, s}}{2} \iint_{\Q} \frac{(u(x)-u(y))(v(x)-v(y))}{|x-y|^{N+2s}} \,dx\,dy\, d\mu(s)\\
&\qquad+\int_\Omega uv\, dx - \lambda\int_\Omega uv\, dx -\int_\Omega f(x, u) v \, dx.
\end{split}
\end{equation}

\begin{proposition}\label{boundedlinking}
Let $\lambda$ be as in~\eqref{successivi} and suppose that~$f$ satisfies~\eqref{AR1}, \eqref{AR2}, \eqref{AR3} and~\eqref{AR4}. Let~$c\in\R$ and let~$u_n$ be a sequence in $\widetilde{\mathcal{H}}_{\alpha, \mu}(\Omega)$ such that
\begin{equation}\label{bdd12}
\lim_{n\to +\infty} \mathcal I(u_n)=c
\end{equation}
and
\begin{equation}\label{bdd22}
\lim_{n\to +\infty} \sup_{\substack{v\in \widetilde{\mathcal{H}}_{\alpha, \mu}(\Omega)\\ \|v\|_{\alpha, \mu}=1}} |\langle \mathcal I'(u_n), v\rangle| =0.
\end{equation}

Then, $u_n$ is bounded in~$\widetilde{\mathcal{H}}_{\alpha, \mu}(\Omega)$.
\end{proposition}

\begin{proof}
Let $k\in (2, \vartheta)$. By~\eqref{bdd12} and~\eqref{bdd22} we infer that there exists a constant $c_1>0$ such that for any $n\in\N$,
\begin{equation}\label{ndkjsc}
k \mathcal I(u_n) -\langle \mathcal I'(u_n), u_n\rangle \le c_1 (k +\|u_n\|_{\alpha, \mu}).
\end{equation}
Let $r$ be as in assumption~\eqref{AR3} and notice that
\begin{equation}\label{lnsdnv}
\int_{\Omega\cap\{|u_n|> r\}} \Big(f(x, u_n) u_n -k F(x, u_n)\Big)\,dx >0.
\end{equation}
Moreover, if $r>0$, then by using Lemma~\ref{SV12} with~$\varepsilon:=1$ we get that
\begin{equation}\label{ler2}
\left|\;\int_{\Omega\cap\{|u_n|\le r\}}  \big(f(x, u_n) u_n -\vartheta F(x, u_n)\big) \,dx\right| \le \left(2r^2 +p\delta(1) r^p + \vartheta r^2 +\vartheta\delta(1) r^p\right)|\Omega|=:C_r.
\end{equation}
Hence, by~\eqref{funlink}, \eqref{diff2},
\eqref{lnsdnv} and~\eqref{ler2}, we infer that
\[
\begin{split}
k\, \mathcal I(u_n) -\langle \mathcal I'(u_n), u_n\rangle &\ge\left(\frac{k}{2}-1\right)\|u_n\|^2_{\alpha, \mu}-\lambda\left(\frac{k}{2}-1\right)\|u_n\|^2_{L^2(\Omega)} \\
&\qquad+\int_\Omega \big(f(x, u_n) u_n -k F(x, u_n)\big) \,dx\\
& \ge\left(\frac{k}{2}-1\right)\|u_n\|^2_{\alpha, \mu}-\lambda\left(\frac{k}{2}-1\right)\|u_n\|^2_{L^2(\Omega)}\\
&\qquad - \left|\;\int_{\Omega\cap\{|u_n|\le r\}}  \Big(f(x, u_n) u_n -\vartheta F(x, u_n)\Big) \,dx\right| +(\vartheta -k)\int_\Omega F(x, u_n) \,dx\\
&\ge\left(\frac{k}{2}-1\right)\|u_n\|^2_{\alpha, \mu}-\lambda\left(\frac{k}{2}-1\right)\|u_n\|^2_{L^2(\Omega)}\\
&\qquad +(\vartheta -k)\int_\Omega F(x, u_n) \,dx- C_r.
\end{split}
\]

Let now~$\widetilde\vartheta$ be as in hypothesis~\eqref{AR4}. We notice that, by using both H\"older and Young inequalities,
for any~$\varepsilon>0$ and for any~$u\in \widetilde{\mathcal{H}}_{\alpha, \mu}(\Omega)$,
\[
\|u\|^2_{L^2(\Omega)}\le \frac{2\,\varepsilon^{\widetilde\vartheta/2}}{\widetilde\vartheta} \|u\|^{\widetilde\vartheta}_{L^{\widetilde\vartheta}(\Omega)} +\frac{\widetilde\vartheta -2}{\widetilde\vartheta\,\varepsilon^{\widetilde\vartheta/{(\widetilde\vartheta -2)}}} |\Omega|.
\]
In light of this, recalling that $\vartheta>k$ and that~\eqref{AR4} holds, we infer that
\[
\begin{split}
k\, \mathcal I(u_n) -\langle \mathcal I'(u_n), u_n\rangle&\ge\left(\frac{k}{2}-1\right)\|u_n\|^2_{\alpha, \mu}-\lambda\left(\frac{k}{2}-1\right)\|u_n\|^2_{L^2(\Omega)}\\
&\qquad +(\vartheta -k) a_3 \|u_n\|^{\widetilde\vartheta}_{L^{\widetilde\vartheta}(\Omega)} +(\vartheta - k)\int_\Omega a_4(x) \,dx -C_r\\
&\ge\left(\frac{k}{2}-1\right)\|u_n\|^2_{\alpha, \mu} + \left((\vartheta-k)a_3-\frac{2\,\lambda}{\widetilde\vartheta}\,\varepsilon^{\widetilde\vartheta/2}\left(\frac{k}{2}-1 \right)\right)\|u_n\|^{\widetilde\vartheta}_{L^{\widetilde\vartheta}(\Omega)} -c_2,
\end{split}
\]
where
\[
c_2:= C_r +\frac{|\Omega|(\widetilde\vartheta-2)}{\widetilde\vartheta\,\varepsilon^{\widetilde\vartheta/(\widetilde\vartheta-2)}} +(\vartheta - k)\int_\Omega a_4(x) \,dx. 
\]

Now, by taking
\[
\varepsilon\in\left(0, \left(\frac{a_3\widetilde\vartheta (\vartheta-k)}{\lambda(k-2)}\right)^{2/\widetilde\vartheta}\right),
\]
we get
\[
k\, \mathcal I(u_n) -\langle \mathcal I'(u_n), u_n\rangle\ge \left(\frac{k}{2}-1\right)\|u_n\|^2_{\alpha, \mu} -c_2.
\]
This, together with~\eqref{ndkjsc}, entails that 
\[
\|u_n\|^2_{\alpha, \mu}\le \frac{2c_1 (k +\|u_n\|_{\alpha, \mu}) +2c_2}{k-2},
\]
namely $u_n$ is bounded in $\widetilde{\mathcal{H}}_{\alpha, \mu}(\Omega)$, as desired.
\end{proof}

\end{document}
